% !PRINTED 313
\weberpage{313}
%<<CONT
\weq{14}{x = y(\lambda+\mu z) = \frac{(\lambda+\mu z)\gamma_2}{z^2+3}}%
\footnote{Cf. \gesperrt{Kiepert}, Auflösung der Gleichung 5ten Grades. Crelles Journal, Bd.\,87, S.\,114.}

are, when $\lambda, \mu$ denote arbitrary parameters which are the same for all
five values $x$.

If one puts this equation into the form

\weq{15}{x^5 - 5ax^3 - 5b\gamma_2 x - bc\gamma_2^2 = 0,}

then, after some calculations making use of the particular values of $x$ for
$\gamma_3=0$ and $\gamma_3=\infty$, the following expressions result for
$a, b, c$:

\weq{16}{\begin{aligned}
a &= 8\lambda^3 - 72\lambda\mu^2 + \gamma_3(\lambda^2\mu - \mu^3),\\
b &= \lambda^4 + 18\lambda^2\mu^2 - 27\mu^4 + \gamma_3\lambda\mu^3,\\
c &= \lambda^5 + 10\lambda^3\mu^2 + 45\lambda\mu^3 + \gamma_3\mu^5. % [sic] German source prints μ^3.
\end{aligned}}

\par\noindent\rule{0.3\linewidth}{0.4pt}\par

One more resolvent of degree 5 is to be formed, making use of the function
$f(\omega)$.

We put:

\weq{17}{\begin{gathered}
u = f(\omega), \quad v_\infty = f(5\omega), \quad v_z = f\left(\frac{\omega+c}{5}\right)\\
c \equiv 0 \pmod{48}, \quad z \equiv c \pmod{5},
\end{gathered}}

and by §\,73 we have between $u$ and $v$ the equation of degree 6:

\weq{18}{u^6 + v^6 - u^5v^5 + 4uv = 0.}

We investigate the influence which the three permutations

\[
(\omega, \omega+2), \quad \left(\omega, -\frac{1}{\omega}\right), \quad
\left(\omega, \frac{\omega-1}{\omega+1}\right)
\]

have on the quantities $v_z$. By these permutations $c$ passes over into
$c_1, c_2, c'$, which by §\,69, (9), (11) and §\,73, (6) are determined by the
congruences

\weq{19}{c_1 \equiv c+2, \quad c_2 \equiv -\frac{1}{c}, \quad c' \equiv \frac{c-1}{c+1} \pmod{5}}

Apart from this change of the index, by §\,73, (4) and (9) the $v$ pass over
into

\weq{20}{e^{-\frac{5\pi i}{12}}v, \quad v, \quad -\frac{\sqrt{2}}{v}.}
